\documentclass[11pt]{article}

\usepackage[margin=1in]{geometry}
\usepackage{booktabs}
\usepackage{amsmath}
\usepackage{graphicx}
\usepackage[hidelinks]{hyperref}
\usepackage{xcolor}
\usepackage{enumitem}

\newcommand{\el}{\texttt{easyloops}}
\newcommand{\code}[1]{\texttt{#1}}

\title{\textbf{Loop Engineering: World-Class Results from Small Local Models\\
via Verification-Anchored Agent Loops}}
\author{Abhimanyu Bhagwati\\
\small\href{https://github.com/AbhimanyuBhagwati/easyloops}{github.com/AbhimanyuBhagwati/easyloops}
\quad\href{https://pypi.org/project/easyloops/}{pypi.org/project/easyloops}}
\date{July 2026}

\begin{document}
\maketitle

\begin{abstract}
Frontier coding agents assume a frontier model. We ask the opposite question: how far can
\emph{small, freely available local models} (2--30B parameters) be pushed by engineering the
\emph{loop} around them rather than scaling the model inside them? We introduce
\textbf{loop engineering}: the discipline of designing verification, context hygiene, recovery
policy, and capability measurement so that unlimited free local inference converts into real
task quality. We present \el{}, an open-source, zero-dependency harness implementing these
principles, and report measurements across six local models on a consumer laptop. Our central
findings: (i)~a 2.5\,GB model (qwen3:4b-instruct) matches models an order of magnitude larger
on tool-driven code repair when anchored to an external verifier, while being the fastest
worker at 151 tokens/s; (ii)~capability must be measured, not assumed---one popular 8B model
proved entirely incapable of tool use under its default runtime, a failure invisible until
probed; and (iii)~the practical quality ceiling of an agent loop is set by its verifier, not
its model: we document failure modes (self-satisfying tests, stale bytecode caches) in which a
loop reports success while doing nothing, and the mechanical defenses that close them.
\end{abstract}

\section{Introduction}

Agent harnesses such as Claude Code demonstrate that a language model equipped with tools and
an execution loop can complete substantial software tasks. These systems, however, presuppose a
highly capable hosted model. Meanwhile, millions of practitioners run small open-weight models
locally---on laptops, without per-token cost, without rate limits, and without sending code to
third parties.

Small models fail differently from large ones. They cannot hold a large plan in context, they
degrade quickly when their context is polluted by their own failed attempts, and their
capabilities vary wildly and undocumented across sizes, families, and runtimes. Naively placing
a small model inside a large-model harness therefore fails. Our thesis is that these failures
are properties of the \emph{loop}, not only of the model, and that they are engineerable:

\begin{quote}
\emph{A small model looping on its own opinion of its output plateaus. A small model looping
against an objective external signal climbs.}
\end{quote}

We call the resulting discipline \textbf{loop engineering} and reduce it to four practices:
verifier design (\S\ref{sec:verify}), capability measurement (\S\ref{sec:bench}), context
hygiene and recovery policy (\S\ref{sec:recovery}), and skills-as-procedure
(\S\ref{sec:skills}). \el{} packages these practices behind a single interactive CLI.

\section{Design}

\subsection{Verification is the ceiling}\label{sec:verify}

Every task produced by the \el{} planner must carry a \code{verify\_cmd}: a shell command that
exits~0 only if the task succeeded (typically running a test file). The harness---never the
model---executes this command to decide pass/fail. Free local compute then behaves like an
optimizer against an objective function: retries are cheap, and each retry receives the failure
output as signal.

The corollary is that verifier defects, not model defects, bound quality. We observed a 4B
worker produce tests that \emph{printed} ``Test failed'' yet exited~0---a verifier that could
never fail, satisfied on the first try. Because a passing verifier exerts no improvement
pressure, the loop happily accepted it. Defenses must be mechanical: \el{} rejects inline
\code{python -c} verifiers at plan validation, ships an auto-attached skill forbidding
print-style tests, and we validate generated tests by \emph{mutation}: deliberately breaking
the implementation and requiring the tests to go red.

\subsection{Measure capability; never assume it}\label{sec:bench}

\code{easyloops bench} probes every installed model with four scored tasks: (1)~\textbf{tools}:
drive a write-file/finish tool protocol; (2)~\textbf{json}: emit a schema-valid task plan;
(3)~\textbf{repair}: fix deliberately buggy code until real tests pass, inside the production
worker loop, with an anti-cheat hash proving the tests were not edited; (4)~\textbf{follow}: an
exact-output instruction probe. Speed and model size are recorded alongside.

Routing is derived from the data by three rules reflecting how each role is used:
the \emph{largest} qualified model plans (planning runs once per goal and gates everything
downstream); the \emph{fastest} top-repairing model works (workers run constantly and
verification catches their mistakes); the \emph{strongest} tool-capable model takes the final
attempt of stuck tasks.

\subsection{Context hygiene and recovery policy}\label{sec:recovery}

Each task attempt starts from \emph{fresh} context: the self-contained task description,
summaries of dependency tasks, a workspace listing, and only the \emph{tail of the previous
failure}---never the prior transcript. Durable state lives on disk in a resumable ledger, not
in any context window. Attempts follow a diversity-then-capability schedule: attempt~1 samples
conservatively (temperature 0.2); attempt~2 repairs in place with the error as input (0.5);
the final attempt \textbf{rolls the workspace back to a clean snapshot} and resamples hot (0.8)
on the escalation model---for small models, a clean retry empirically beats digging deeper into
a broken repair. A task that exhausts its attempts is returned to the planner with its failure
output and \emph{split} into smaller subtasks, the dependency DAG rewiring automatically.

\subsection{Skills as written procedure}\label{sec:skills}

Small models benefit disproportionately from explicit procedure. \el{} skills are Markdown
files carrying a description and trigger keywords; they attach to tasks either by planner
choice or---because planners proved unreliable at opting in---mechanically, by keyword match
against the task text. Their bodies are injected into the worker's system prompt.

\section{Evaluation}

\subsection{Capability probes}

Table~\ref{tab:bench} reports bench results on an Apple M3~Max (36\,GB unified memory,
Ollama 0.30.5). Two results stand out. First, \textbf{qwen3:4b-instruct}---a 2.5\,GB
model---achieves perfect scores on every probe while being the fastest model measured; it
repaired the buggy module in 5.3\,s, faster than models $7\times$ its size. Second,
\textbf{llama3:8b scores zero on tool use}: its runtime rejects tool schemas outright. A user
pairing that model with any agent harness would fail opaquely mid-task; the probe surfaces the
incompatibility in seconds.

\begin{table}[h]
\centering
\caption{Capability probes across six local models (scores in $[0,1]$).}
\label{tab:bench}
\begin{tabular}{lccccr}
\toprule
Model & Tools & Plan JSON & Repair & Follow & tok/s \\
\midrule
qwen3:4b-instruct  & 1.0 & 1.0 & 1.0 (5.3\,s)  & 1.0 & 151 \\
qwen3:30b-instruct & 1.0 & 1.0 & 1.0 (7.2\,s)  & 1.0 & 128 \\
gemma4             & 1.0 & 1.0 & 1.0 (17.2\,s) & 1.0 & 103 \\
qwen2.5:14b        & 1.0 & 1.0 & 1.0 (26.6\,s) & 1.0 & 54  \\
qwen2.5:7b         & 1.0 & 1.0 & 0.7 (2 attempts) & 1.0 & 103 \\
llama3:8b          & \textbf{0.0} & 1.0 & 0.0 & 1.0 & 94 \\
\bottomrule
\end{tabular}
\end{table}

Derived routing: planner and escalation \code{qwen3:30b-instruct}; worker and utility
\code{qwen3:4b-instruct}.

\subsection{Case studies}

\textbf{End-to-end builds.} Given natural-language goals (a JSON-backed expense tracker with an
\code{argparse} CLI and tests; a timestamped notes module with tests), the harness planned 2--5
task DAGs and completed them with the 4B worker, most tasks passing on the first attempt. All
artifacts were verified independently of the harness's own report.

\textbf{The full safety net, organically.} In the notes-module build, the 4B worker failed the
test-writing task twice; the harness rolled the workspace back to its snapshot, escalated the
final attempt to the 30B at temperature~0.8, and passed. The resulting test suite survived
mutation testing (three deliberate implementation bugs, all caught).

\textbf{Verifier blindness and its repair.} The expense-tracker run initially yielded the
print-style non-failing tests of \S\ref{sec:verify}. Re-routing only the test-writing task to
the 30B---with the anchor file removed and the anti-pattern skill attached---produced 25 real
assertions; mutation testing confirmed them. The general lesson: tasks whose output \emph{is}
the verifier are precisely the tasks that need the strongest model.

\textbf{A systems pitfall: stale bytecode.} During evaluation we observed verify commands
executing \emph{stale code}: CPython invalidates cached bytecode by (mtime-seconds, size), so a
same-size rewrite within one second---routine in fast agent loops---silently reuses the old
compile. On macOS the system Python further relocates caches to a hidden per-user directory.
\el{} now assigns each workspace a private \code{PYTHONPYCACHEPREFIX}, wiped before every
command, making verify results reflect the files actually on disk. We flag this as a hazard for
any agent framework that executes rapidly rewritten Python.

\subsection{Robustness}

Real home directories are hostile: we encountered a symlink cycle (created by an unrelated
tool) that drove a naive directory walker past the OS path-length limit and crashed the
session. The \el{} walker now refuses symlinks, caps depth and file count, and treats
unstatable entries as skippable; the interactive CLI additionally refuses to adopt \code{\$HOME}
or a filesystem root as an agent workspace. Both behaviors are regression-tested, including
against the original pathological directory.

\section{Limitations}

The approach inherits the limits of its verifiers: tasks with machine-checkable success
(code-with-tests, data transformation, format conversion) climb under iteration; open-ended
taste tasks do not, because nothing objective tells the loop it is wrong. Our probes saturate
on capable models---they separate \emph{can} from \emph{cannot} sharply but not \emph{good}
from \emph{great}. The workspace jail is a convention, not a sandbox. Single-machine memory
(one resident large model) serializes execution; we do not yet exploit parallel sampling.

\section{Future Work}

\textbf{Managed-cloud backends.} \el{} speaks two protocols: native Ollama and OpenAI-compatible
chat completions. Because AWS Bedrock and Google Vertex~AI both expose OpenAI-compatible
endpoints, they are reachable today via \code{-{}-backend openai}; first-class support requires
native credential handling (AWS SigV4/IAM; Google OAuth tokens with hourly expiry) and model
discovery, which we plan as dedicated backend adapters. The interesting research question is
economic: with a metered cloud model in the escalation slot and free local models in the worker
slot, the routing rules of \S\ref{sec:bench} become a cost optimizer---spend cents only where
verification proves the free tier insufficient.

\textbf{Also planned:} best-of-$N$ verified sampling under runtime parallelism; embedding-based
file retrieval for large workspaces; a built-in mutation check after every test-writing task;
harder, discriminative probes; and task-type routing (quality-sensitive tasks straight to the
strongest model).

\section{Conclusion}

Loop engineering reframes the small-model question. The scarce resource is not intelligence per
token but \emph{verifiable feedback per attempt}. Given an objective verifier, honest
measurement of what each local model can do, ruthless context hygiene, and a recovery ladder of
rollback, escalation, and re-planning, a 2.5\,GB model on a laptop does real software work---and
knows when to hand the keyboard to a bigger sibling. All code, benchmarks, and the harness
itself are open source: \code{pip install easyloops}.

\end{document}
